% Options for packages loaded elsewhere
% Options for packages loaded elsewhere
\PassOptionsToPackage{unicode}{hyperref}
\PassOptionsToPackage{hyphens}{url}
\PassOptionsToPackage{dvipsnames,svgnames,x11names}{xcolor}
%
\documentclass[
  french,
  letterpaper,
  DIV=11,
  numbers=noendperiod]{scrartcl}
\usepackage{xcolor}
\usepackage{amsmath,amssymb}
\setcounter{secnumdepth}{5}
\usepackage{iftex}
\ifPDFTeX
  \usepackage[T1]{fontenc}
  \usepackage[utf8]{inputenc}
  \usepackage{textcomp} % provide euro and other symbols
\else % if luatex or xetex
  \usepackage{unicode-math} % this also loads fontspec
  \defaultfontfeatures{Scale=MatchLowercase}
  \defaultfontfeatures[\rmfamily]{Ligatures=TeX,Scale=1}
\fi
\usepackage{lmodern}
\ifPDFTeX\else
  % xetex/luatex font selection
\fi
% Use upquote if available, for straight quotes in verbatim environments
\IfFileExists{upquote.sty}{\usepackage{upquote}}{}
\IfFileExists{microtype.sty}{% use microtype if available
  \usepackage[]{microtype}
  \UseMicrotypeSet[protrusion]{basicmath} % disable protrusion for tt fonts
}{}
\makeatletter
\@ifundefined{KOMAClassName}{% if non-KOMA class
  \IfFileExists{parskip.sty}{%
    \usepackage{parskip}
  }{% else
    \setlength{\parindent}{0pt}
    \setlength{\parskip}{6pt plus 2pt minus 1pt}}
}{% if KOMA class
  \KOMAoptions{parskip=half}}
\makeatother
% Make \paragraph and \subparagraph free-standing
\makeatletter
\ifx\paragraph\undefined\else
  \let\oldparagraph\paragraph
  \renewcommand{\paragraph}{
    \@ifstar
      \xxxParagraphStar
      \xxxParagraphNoStar
  }
  \newcommand{\xxxParagraphStar}[1]{\oldparagraph*{#1}\mbox{}}
  \newcommand{\xxxParagraphNoStar}[1]{\oldparagraph{#1}\mbox{}}
\fi
\ifx\subparagraph\undefined\else
  \let\oldsubparagraph\subparagraph
  \renewcommand{\subparagraph}{
    \@ifstar
      \xxxSubParagraphStar
      \xxxSubParagraphNoStar
  }
  \newcommand{\xxxSubParagraphStar}[1]{\oldsubparagraph*{#1}\mbox{}}
  \newcommand{\xxxSubParagraphNoStar}[1]{\oldsubparagraph{#1}\mbox{}}
\fi
\makeatother

\usepackage{color}
\usepackage{fancyvrb}
\newcommand{\VerbBar}{|}
\newcommand{\VERB}{\Verb[commandchars=\\\{\}]}
\DefineVerbatimEnvironment{Highlighting}{Verbatim}{commandchars=\\\{\}}
% Add ',fontsize=\small' for more characters per line
\usepackage{framed}
\definecolor{shadecolor}{RGB}{241,243,245}
\newenvironment{Shaded}{\begin{snugshade}}{\end{snugshade}}
\newcommand{\AlertTok}[1]{\textcolor[rgb]{0.68,0.00,0.00}{#1}}
\newcommand{\AnnotationTok}[1]{\textcolor[rgb]{0.37,0.37,0.37}{#1}}
\newcommand{\AttributeTok}[1]{\textcolor[rgb]{0.40,0.45,0.13}{#1}}
\newcommand{\BaseNTok}[1]{\textcolor[rgb]{0.68,0.00,0.00}{#1}}
\newcommand{\BuiltInTok}[1]{\textcolor[rgb]{0.00,0.23,0.31}{#1}}
\newcommand{\CharTok}[1]{\textcolor[rgb]{0.13,0.47,0.30}{#1}}
\newcommand{\CommentTok}[1]{\textcolor[rgb]{0.37,0.37,0.37}{#1}}
\newcommand{\CommentVarTok}[1]{\textcolor[rgb]{0.37,0.37,0.37}{\textit{#1}}}
\newcommand{\ConstantTok}[1]{\textcolor[rgb]{0.56,0.35,0.01}{#1}}
\newcommand{\ControlFlowTok}[1]{\textcolor[rgb]{0.00,0.23,0.31}{\textbf{#1}}}
\newcommand{\DataTypeTok}[1]{\textcolor[rgb]{0.68,0.00,0.00}{#1}}
\newcommand{\DecValTok}[1]{\textcolor[rgb]{0.68,0.00,0.00}{#1}}
\newcommand{\DocumentationTok}[1]{\textcolor[rgb]{0.37,0.37,0.37}{\textit{#1}}}
\newcommand{\ErrorTok}[1]{\textcolor[rgb]{0.68,0.00,0.00}{#1}}
\newcommand{\ExtensionTok}[1]{\textcolor[rgb]{0.00,0.23,0.31}{#1}}
\newcommand{\FloatTok}[1]{\textcolor[rgb]{0.68,0.00,0.00}{#1}}
\newcommand{\FunctionTok}[1]{\textcolor[rgb]{0.28,0.35,0.67}{#1}}
\newcommand{\ImportTok}[1]{\textcolor[rgb]{0.00,0.46,0.62}{#1}}
\newcommand{\InformationTok}[1]{\textcolor[rgb]{0.37,0.37,0.37}{#1}}
\newcommand{\KeywordTok}[1]{\textcolor[rgb]{0.00,0.23,0.31}{\textbf{#1}}}
\newcommand{\NormalTok}[1]{\textcolor[rgb]{0.00,0.23,0.31}{#1}}
\newcommand{\OperatorTok}[1]{\textcolor[rgb]{0.37,0.37,0.37}{#1}}
\newcommand{\OtherTok}[1]{\textcolor[rgb]{0.00,0.23,0.31}{#1}}
\newcommand{\PreprocessorTok}[1]{\textcolor[rgb]{0.68,0.00,0.00}{#1}}
\newcommand{\RegionMarkerTok}[1]{\textcolor[rgb]{0.00,0.23,0.31}{#1}}
\newcommand{\SpecialCharTok}[1]{\textcolor[rgb]{0.37,0.37,0.37}{#1}}
\newcommand{\SpecialStringTok}[1]{\textcolor[rgb]{0.13,0.47,0.30}{#1}}
\newcommand{\StringTok}[1]{\textcolor[rgb]{0.13,0.47,0.30}{#1}}
\newcommand{\VariableTok}[1]{\textcolor[rgb]{0.07,0.07,0.07}{#1}}
\newcommand{\VerbatimStringTok}[1]{\textcolor[rgb]{0.13,0.47,0.30}{#1}}
\newcommand{\WarningTok}[1]{\textcolor[rgb]{0.37,0.37,0.37}{\textit{#1}}}

\usepackage{longtable,booktabs,array}
\usepackage{calc} % for calculating minipage widths
% Correct order of tables after \paragraph or \subparagraph
\usepackage{etoolbox}
\makeatletter
\patchcmd\longtable{\par}{\if@noskipsec\mbox{}\fi\par}{}{}
\makeatother
% Allow footnotes in longtable head/foot
\IfFileExists{footnotehyper.sty}{\usepackage{footnotehyper}}{\usepackage{footnote}}
\makesavenoteenv{longtable}
\usepackage{graphicx}
\makeatletter
\newsavebox\pandoc@box
\newcommand*\pandocbounded[1]{% scales image to fit in text height/width
  \sbox\pandoc@box{#1}%
  \Gscale@div\@tempa{\textheight}{\dimexpr\ht\pandoc@box+\dp\pandoc@box\relax}%
  \Gscale@div\@tempb{\linewidth}{\wd\pandoc@box}%
  \ifdim\@tempb\p@<\@tempa\p@\let\@tempa\@tempb\fi% select the smaller of both
  \ifdim\@tempa\p@<\p@\scalebox{\@tempa}{\usebox\pandoc@box}%
  \else\usebox{\pandoc@box}%
  \fi%
}
% Set default figure placement to htbp
\def\fps@figure{htbp}
\makeatother



\ifLuaTeX
\usepackage[bidi=basic]{babel}
\else
\usepackage[bidi=default]{babel}
\fi
% get rid of language-specific shorthands (see #6817):
\let\LanguageShortHands\languageshorthands
\def\languageshorthands#1{}


\setlength{\emergencystretch}{3em} % prevent overfull lines

\providecommand{\tightlist}{%
  \setlength{\itemsep}{0pt}\setlength{\parskip}{0pt}}



 


\KOMAoption{captions}{tableheading}
\makeatletter
\@ifpackageloaded{tcolorbox}{}{\usepackage[skins,breakable]{tcolorbox}}
\@ifpackageloaded{fontawesome5}{}{\usepackage{fontawesome5}}
\definecolor{quarto-callout-color}{HTML}{909090}
\definecolor{quarto-callout-note-color}{HTML}{0758E5}
\definecolor{quarto-callout-important-color}{HTML}{CC1914}
\definecolor{quarto-callout-warning-color}{HTML}{EB9113}
\definecolor{quarto-callout-tip-color}{HTML}{00A047}
\definecolor{quarto-callout-caution-color}{HTML}{FC5300}
\definecolor{quarto-callout-color-frame}{HTML}{acacac}
\definecolor{quarto-callout-note-color-frame}{HTML}{4582ec}
\definecolor{quarto-callout-important-color-frame}{HTML}{d9534f}
\definecolor{quarto-callout-warning-color-frame}{HTML}{f0ad4e}
\definecolor{quarto-callout-tip-color-frame}{HTML}{02b875}
\definecolor{quarto-callout-caution-color-frame}{HTML}{fd7e14}
\makeatother
\makeatletter
\@ifpackageloaded{caption}{}{\usepackage{caption}}
\AtBeginDocument{%
\ifdefined\contentsname
  \renewcommand*\contentsname{Table des matières}
\else
  \newcommand\contentsname{Table des matières}
\fi
\ifdefined\listfigurename
  \renewcommand*\listfigurename{Liste des Figures}
\else
  \newcommand\listfigurename{Liste des Figures}
\fi
\ifdefined\listtablename
  \renewcommand*\listtablename{Liste des Tables}
\else
  \newcommand\listtablename{Liste des Tables}
\fi
\ifdefined\figurename
  \renewcommand*\figurename{Figure}
\else
  \newcommand\figurename{Figure}
\fi
\ifdefined\tablename
  \renewcommand*\tablename{Table}
\else
  \newcommand\tablename{Table}
\fi
}
\@ifpackageloaded{float}{}{\usepackage{float}}
\floatstyle{ruled}
\@ifundefined{c@chapter}{\newfloat{codelisting}{h}{lop}}{\newfloat{codelisting}{h}{lop}[chapter]}
\floatname{codelisting}{Listing}
\newcommand*\listoflistings{\listof{codelisting}{Liste des Listings}}
\makeatother
\makeatletter
\makeatother
\makeatletter
\@ifpackageloaded{caption}{}{\usepackage{caption}}
\@ifpackageloaded{subcaption}{}{\usepackage{subcaption}}
\makeatother
\usepackage{bookmark}
\IfFileExists{xurl.sty}{\usepackage{xurl}}{} % add URL line breaks if available
\urlstyle{same}
\hypersetup{
  pdftitle={Exercice : Intégration de l'API eBay},
  pdflang={fr},
  colorlinks=true,
  linkcolor={blue},
  filecolor={Maroon},
  citecolor={Blue},
  urlcolor={Blue},
  pdfcreator={LaTeX via pandoc}}


\title{Exercice : Intégration de l'API eBay}
\usepackage{etoolbox}
\makeatletter
\providecommand{\subtitle}[1]{% add subtitle to \maketitle
  \apptocmd{\@title}{\par {\large #1 \par}}{}{}
}
\makeatother
\subtitle{Client Credentials Grant --- Browse API}
\author{}
\date{}
\begin{document}
\maketitle

\renewcommand*\contentsname{Sommaire}
{
\hypersetup{linkcolor=}
\setcounter{tocdepth}{3}
\tableofcontents
}

\section{Contexte}\label{contexte}

Dans cet exercice, vous allez implémenter un client Python pour
interroger l'\textbf{API eBay}. Vous utiliserez le flux \textbf{OAuth
Client Credentials Grant} pour obtenir un token d'accès, puis
effectuerez des recherches d'articles via la \textbf{Browse API}.

Ce type d'authentification (sans utilisateur connecté) est adapté aux
accès en lecture seule aux données publiques eBay.

\begin{tcolorbox}[enhanced jigsaw, toptitle=1mm, arc=.35mm, colbacktitle=quarto-callout-note-color!10!white, title=\textcolor{quarto-callout-note-color}{\faInfo}\hspace{0.5em}{Ressources}, titlerule=0mm, toprule=.15mm, rightrule=.15mm, leftrule=.75mm, breakable, left=2mm, opacityback=0, coltitle=black, bottomrule=.15mm, colframe=quarto-callout-note-color-frame, bottomtitle=1mm, colback=white, opacitybacktitle=0.6]

\begin{itemize}
\tightlist
\item
  \href{https://developer.ebay.com/api-docs/static/oauth-client-credentials-grant.html}{OAuth
  Client Credentials Grant}
\item
  \href{https://developer.ebay.com/my/api_test_tool?index=0&env=production}{Outil
  de test API eBay}
\item
  \href{https://developer.ebay.com/api-docs/buy/browse/resources/item_summary/methods/search}{Browse
  API --- Item Summary Search}
\end{itemize}

\end{tcolorbox}

\begin{center}\rule{0.5\linewidth}{0.5pt}\end{center}

\section{Prérequis}\label{pruxe9requis}

Avant de commencer, assurez-vous d'avoir :

\begin{itemize}
\tightlist
\item
  Un compte développeur eBay avec une application enregistrée
  (\textbf{App ID} et \textbf{Client Secret})
\item
  Python 3.8+ avec les bibliothèques \texttt{requests} et
  \texttt{python-dotenv}
\item
  Un fichier \texttt{.env} à la racine du projet contenant vos
  identifiants
\end{itemize}

\begin{Shaded}
\begin{Highlighting}[]
\CommentTok{\# .env}
\VariableTok{app\_id}\OperatorTok{=}\NormalTok{VotreAppID}
\VariableTok{dev\_id}\OperatorTok{=}\NormalTok{VotreDevID}
\VariableTok{client\_secret}\OperatorTok{=}\NormalTok{VotreClientSecret}
\end{Highlighting}
\end{Shaded}

\begin{tcolorbox}[enhanced jigsaw, toptitle=1mm, arc=.35mm, colbacktitle=quarto-callout-warning-color!10!white, title=\textcolor{quarto-callout-warning-color}{\faExclamationTriangle}\hspace{0.5em}{Ne commitez jamais votre fichier .env}, titlerule=0mm, toprule=.15mm, rightrule=.15mm, leftrule=.75mm, breakable, left=2mm, opacityback=0, coltitle=black, bottomrule=.15mm, colframe=quarto-callout-warning-color-frame, bottomtitle=1mm, colback=white, opacitybacktitle=0.6]

Ajoutez \texttt{.env} à votre \texttt{.gitignore} pour ne pas exposer
vos credentials.

\end{tcolorbox}

\begin{center}\rule{0.5\linewidth}{0.5pt}\end{center}

\section{Objectifs}\label{objectifs}

Vous devrez implémenter une classe \texttt{EbayApi} qui :

\begin{enumerate}
\def\labelenumi{\arabic{enumi}.}
\tightlist
\item
  Charge les credentials depuis les variables d'environnement au moment
  de l'instanciation
\item
  Gère un token OAuth avec renouvellement automatique à l'expiration
\item
  Expose les fonctionnalités de recherche à ses utilisateurs
\end{enumerate}

Réfléchissez à la décomposition en méthodes : quelles responsabilités
méritent d'être isolées ? Lesquelles font partie de l'interface
publique, lesquelles sont des détails d'implémentation ?

\begin{center}\rule{0.5\linewidth}{0.5pt}\end{center}

\section{Partie 1 --- Authentification
OAuth}\label{partie-1-authentification-oauth}

\subsection{1.1 Obtenir un token
d'accès}\label{obtenir-un-token-daccuxe8s}

L'API eBay utilise le flux \textbf{Client Credentials Grant}. Vous devez
envoyer une requête \texttt{POST} à l'endpoint :

\begin{verbatim}
POST https://api.ebay.com/identity/v1/oauth2/token
\end{verbatim}

\begin{longtable}[]{@{}ll@{}}
\toprule\noalign{}
Element & Valeur \\
\midrule\noalign{}
\endhead
\bottomrule\noalign{}
\endlastfoot
Header \texttt{Authorization} &
\texttt{Basic\ \textless{}base64(app\_id:client\_secret)\textgreater{}} \\
Header \texttt{Content-Type} &
\texttt{application/x-www-form-urlencoded} \\
Body \texttt{grant\_type} & \texttt{client\_credentials} \\
Body \texttt{scope} & \texttt{https://api.ebay.com/oauth/api\_scope} \\
\end{longtable}

\begin{tcolorbox}[enhanced jigsaw, toptitle=1mm, arc=.35mm, colbacktitle=quarto-callout-tip-color!10!white, title=\textcolor{quarto-callout-tip-color}{\faLightbulb}\hspace{0.5em}{Encodage Base64}, titlerule=0mm, toprule=.15mm, rightrule=.15mm, leftrule=.75mm, breakable, left=2mm, opacityback=0, coltitle=black, bottomrule=.15mm, colframe=quarto-callout-tip-color-frame, bottomtitle=1mm, colback=white, opacitybacktitle=0.6]

En Python, pour encoder les credentials :

\begin{Shaded}
\begin{Highlighting}[]
\ImportTok{import}\NormalTok{ base64}
\NormalTok{credentials }\OperatorTok{=} \SpecialStringTok{f"}\SpecialCharTok{\{}\NormalTok{app\_id}\SpecialCharTok{\}}\SpecialStringTok{:}\SpecialCharTok{\{}\NormalTok{client\_secret}\SpecialCharTok{\}}\SpecialStringTok{"}
\NormalTok{encoded }\OperatorTok{=}\NormalTok{ base64.b64encode(credentials.encode()).decode()}
\end{Highlighting}
\end{Shaded}

\end{tcolorbox}

\subsection{1.2 Gestion de l'expiration}\label{gestion-de-lexpiration}

La réponse contient un champ \texttt{expires\_in} (en secondes,
généralement \texttt{7200}). Le token ne doit être renouvelé que
lorsqu'il est expiré, et ce renouvellement doit être
\textbf{transparent} pour le code appelant.

Réfléchissez à comment Python vous permet d'exposer cela proprement :
comment rendre l'accès au token aussi simple que l'accès à un attribut,
tout en cachant la logique de refresh ?

\begin{center}\rule{0.5\linewidth}{0.5pt}\end{center}

\section{Partie 2 --- Requêtes à la Browse
API}\label{partie-2-requuxeates-uxe0-la-browse-api}

\subsection{2.1 Recherche par mot-clé}\label{recherche-par-mot-cluxe9}

Implémentez une méthode permettant de rechercher des articles à partir
d'un mot-clé :

\begin{verbatim}
GET https://api.ebay.com/buy/browse/v1/item_summary/search?q=<query>
\end{verbatim}

\textbf{Headers obligatoires pour toutes les requêtes :}

\begin{Shaded}
\begin{Highlighting}[]
\NormalTok{\{}
    \StringTok{"Authorization"}\NormalTok{: }\SpecialStringTok{f"Bearer }\SpecialCharTok{\{}\NormalTok{access\_token}\SpecialCharTok{\}}\SpecialStringTok{"}\NormalTok{,}
    \StringTok{"X{-}EBAY{-}C{-}MARKETPLACE{-}ID"}\NormalTok{: }\StringTok{"EBAY\_FR"}\NormalTok{,}
\NormalTok{\}}
\end{Highlighting}
\end{Shaded}

Toutes vos requêtes GET partagent la même logique (headers, gestion des
erreurs). Pensez à factoriser ce comportement commun plutôt que de le
dupliquer.

\subsection{2.2 Récupération d'un article par
URL}\label{ruxe9cupuxe9ration-dun-article-par-url}

Implémentez une méthode permettant de récupérer un article à partir de
son URL complète. Exemple :

\begin{verbatim}
https://api.ebay.com/buy/browse/v1/item/v1%7C358311088969%7C626586956773
\end{verbatim}

\begin{center}\rule{0.5\linewidth}{0.5pt}\end{center}

\section{Partie 3 --- Gestion des erreurs et
logs}\label{partie-3-gestion-des-erreurs-et-logs}

Votre implémentation doit gérer les cas suivants :

\begin{longtable}[]{@{}
  >{\raggedright\arraybackslash}p{(\linewidth - 2\tabcolsep) * \real{0.1923}}
  >{\raggedright\arraybackslash}p{(\linewidth - 2\tabcolsep) * \real{0.8077}}@{}}
\toprule\noalign{}
\begin{minipage}[b]{\linewidth}\raggedright
Cas
\end{minipage} & \begin{minipage}[b]{\linewidth}\raggedright
Comportement attendu
\end{minipage} \\
\midrule\noalign{}
\endhead
\bottomrule\noalign{}
\endlastfoot
Credentials manquants dans \texttt{.env} & Lever une \texttt{ValueError}
explicite \\
Echec token (status != 200) & Lever une \texttt{Exception} avec status +
message \\
Echec requête API (status != 200) & Lever une \texttt{Exception} avec
status + corps \\
Opérations importantes & Utiliser le module \texttt{logging} --- aucun
\texttt{print} \\
\end{longtable}

\begin{center}\rule{0.5\linewidth}{0.5pt}\end{center}

\section{Livrables attendus}\label{livrables-attendus}

Un fichier \texttt{ebay\_api.py} contenant la classe \texttt{EbayApi}
avec toutes les fonctionnalités décrites, ainsi qu'un bloc
\texttt{main()} qui instancie la classe, effectue une recherche et
sauvegarde le résultat dans un fichier JSON.

\begin{center}\rule{0.5\linewidth}{0.5pt}\end{center}

\begin{tcolorbox}[enhanced jigsaw, toptitle=1mm, arc=.35mm, colbacktitle=quarto-callout-tip-color!10!white, title=\textcolor{quarto-callout-tip-color}{\faLightbulb}\hspace{0.5em}{Conseil}, titlerule=0mm, toprule=.15mm, rightrule=.15mm, leftrule=.75mm, breakable, left=2mm, opacityback=0, coltitle=black, bottomrule=.15mm, colframe=quarto-callout-tip-color-frame, bottomtitle=1mm, colback=white, opacitybacktitle=0.6]

Utilisez
l'\href{https://developer.ebay.com/my/api_test_tool?index=0&env=production}{outil
de test eBay} pour vérifier vos credentials et visualiser les réponses
avant de commencer à coder.

\end{tcolorbox}




\end{document}
